% GENERATED FILE. Edit canonical lessons, then run npm run generate:books.
% canonical-source-hash: fnv1a64:f93806ee8c502d7c
% canonical-lessons: UR-C288-chiragh, UR-C288-shola, UR-C288-dahakna, UR-C288-raushni, UR-C288-tapish

\chapter{The Lamp That Fell}
\label{ch:ur-the-lamp-that-fell}
\hlchaptermodality{288}

\begin{chapteropening}
\textbf{By the end of this chapter:} \emph{I can describe a fire at its height.}

Complete the last lesson of chapter 288: I can describe a fire at its height.
\end{chapteropening}

\section[\texorpdfstring{\ur{چراغ} (chirāgh)}{چراغ (chirāgh)}]{\ur{چراغ} (chirāgh) — an oil lamp}

\label{lesson:UR-C288-chiragh}



\begin{quote}
Before the new one: say the Urdu for to flare up, then the Urdu for to be on fire.
\end{quote}

\subsection*{You'll want to know: \ur{چراغ}}
\textbf{\ur{چراغ}} — \emph{chirāgh} — \textquotedblleft{}an oil lamp\textquotedblright{}.

\begin{quote}
\ur{رات} \ur{کو} \ur{چراغ} \ur{گر} \ur{گیا۔}
\end{quote}

\begin{quote}
\emph{rāt ko chirāgh gir gayā.}
\end{quote}

\begin{quote}
The lamp fell over in the night.
\end{quote}

Masculine.

\begin{cousinweb}
\textbf{\ur{چراغ}} (\emph{chirāgh}) is usually traced to Persian \emph{chirāgh}, \textquotedblleft{}a lamp\textquotedblright{}. It is the small clay lamp with oil and a wick.
\end{cousinweb}

\begin{grammarlens}[title={What you've built}]
One more word for telling what happened.
\end{grammarlens}

\subsection*{Your turn}
\begin{itemize}
\raggedright
  \item \emph{Say it:} \emph{chirāgh}
  \item \emph{Say it:} \emph{chirāgh}, once more
  \item \emph{Say it:} \emph{rāt ko chirāgh gir gayā.} — The lamp fell over in the night.
  \item \emph{Recall:} say the Urdu for to flare up, then the Urdu for to be on fire, then say \emph{chirāgh} again
\end{itemize}

\begin{tcolorbox}[breakable,colback=teal!4,colframe=teal!35!black,title={Before you move on}]
What does \ur{چراغ} mean? (\textquotedblleft{}An oil lamp\textquotedblright{}.) Say it once more.
\end{tcolorbox}
\section[\texorpdfstring{\ur{شعلہ} (sholā)}{شعلہ (sholā)}]{\ur{شعلہ} (sholā) — a flame}

\label{lesson:UR-C288-shola}



\begin{quote}
Before the new one: say the Urdu for to be on fire, then the Urdu for an oil lamp.
\end{quote}

\subsection*{You'll want to know: \ur{شعلہ}}
\textbf{\ur{شعلہ}} — \emph{sholā} — \textquotedblleft{}a flame\textquotedblright{}.

\begin{quote}
\ur{شعلے} \ur{چھت} \ur{تک} \ur{پہنچے۔}
\end{quote}

\begin{quote}
\emph{shole chhat tak pahuṅche.}
\end{quote}

\begin{quote}
The flames reached the roof.
\end{quote}

Masculine.

\begin{cousinweb}
\textbf{\ur{شعلہ}} (\emph{sholā}) is usually traced, through Persian, to Arabic \emph{shuʿla}, \textquotedblleft{}a flame\textquotedblright{}. The plural is \textbf{\ur{شعلے}}.
\end{cousinweb}

\begin{grammarlens}[title={What you've built}]
One more word for telling what happened.
\end{grammarlens}

\subsection*{Your turn}
\begin{itemize}
\raggedright
  \item \emph{Say it:} \emph{sholā}
  \item \emph{Say it:} \emph{sholā}, once more
  \item \emph{Say it:} \emph{shole chhat tak pahuṅche.} — The flames reached the roof.
  \item \emph{Recall:} say the Urdu for to be on fire, then the Urdu for an oil lamp, then say \emph{sholā} again
\end{itemize}

\begin{tcolorbox}[breakable,colback=teal!4,colframe=teal!35!black,title={Before you move on}]
What does \ur{شعلہ} mean? (\textquotedblleft{}A flame\textquotedblright{}.) Say it once more.
\end{tcolorbox}
\section[\texorpdfstring{\ur{دہکنا} (dahaknā)}{دہکنا (dahaknā)}]{\ur{دہکنا} (dahaknā) — to glow red-hot}

\label{lesson:UR-C288-dahakna}



\begin{quote}
Before the new one: say the Urdu for an oil lamp, then the Urdu for a flame.
\end{quote}

\subsection*{You'll want to know: \ur{دہکنا}}
\textbf{\ur{دہکنا}} — \emph{dahaknā} — \textquotedblleft{}to glow red-hot\textquotedblright{}.

\begin{quote}
\ur{انگارے} \ur{دہک} \ur{رہے} \ur{تھے۔}
\end{quote}

\begin{quote}
\emph{angāre dahak rahe the.}
\end{quote}

\begin{quote}
The embers were glowing red.
\end{quote}

Past (it …): \ur{وہ} \ur{انگارا} \ur{دہکا}, \emph{voh angārā dahkā}.

\begin{cousinweb}
\textbf{\ur{دہکنا}} (\emph{dahaknā}) is inherited: CDIAL lists the Hindi cousin \emph{dahaknā}, \textquotedblleft{}to blaze up\textquotedblright{}, under Sanskrit \emph{dáhati}, \textquotedblleft{}burns\textquotedblright{} (CDIAL 6245). It is fire without flame, glowing.
\end{cousinweb}

\begin{grammarlens}[title={What you've built}]
One more word for telling what happened.
\end{grammarlens}

\subsection*{Your turn}
\begin{itemize}
\raggedright
  \item \emph{Say it:} \emph{dahaknā}
  \item \emph{Say it:} \emph{dahaknā}, once more
  \item \emph{Say it:} \emph{angāre dahak rahe the.} — The embers were glowing red.
  \item \emph{Recall:} say the Urdu for an oil lamp, then the Urdu for a flame, then say \emph{dahaknā} again
\end{itemize}

\begin{tcolorbox}[breakable,colback=teal!4,colframe=teal!35!black,title={Before you move on}]
What does \ur{دہکنا} mean? (\textquotedblleft{}To glow red-hot\textquotedblright{}.) Say it once more.
\end{tcolorbox}
\section[\texorpdfstring{\ur{روشنی} (raushnī)}{روشنی (raushnī)}]{\ur{روشنی} (raushnī) — light; brightness}

\label{lesson:UR-C288-raushni}



\begin{quote}
Before the new one: say the Urdu for a flame, then the Urdu for to glow red-hot.
\end{quote}

\subsection*{You'll want to know: \ur{روشنی}}
\textbf{\ur{روشنی}} — \emph{raushnī} — \textquotedblleft{}light; brightness\textquotedblright{}.

\begin{quote}
\ur{آگ} \ur{کی} \ur{روشنی} \ur{دور} \ur{تک} \ur{تھی۔}
\end{quote}

\begin{quote}
\emph{āg kī raushnī dūr tak thī.}
\end{quote}

\begin{quote}
The light of the fire could be seen far off.
\end{quote}

Feminine.

\begin{cousinweb}
\textbf{\ur{روشنی}} (\emph{raushnī}) is usually traced to Persian \emph{raushanī}, \textquotedblleft{}brightness\textquotedblright{}, from \emph{raushan}, \textquotedblleft{}bright\textquotedblright{}. It is feminine: \textbf{\ur{روشنی} \ur{تھی}}.
\end{cousinweb}

\begin{grammarlens}[title={What you've built}]
One more word for telling what happened.
\end{grammarlens}

\subsection*{Your turn}
\begin{itemize}
\raggedright
  \item \emph{Say it:} \emph{raushnī}
  \item \emph{Say it:} \emph{raushnī}, once more
  \item \emph{Say it:} \emph{āg kī raushnī dūr tak thī.} — The light of the fire could be seen far off.
  \item \emph{Recall:} say the Urdu for a flame, then the Urdu for to glow red-hot, then say \emph{raushnī} again
\end{itemize}

\begin{tcolorbox}[breakable,colback=teal!4,colframe=teal!35!black,title={Before you move on}]
What does \ur{روشنی} mean? (\textquotedblleft{}Light; brightness\textquotedblright{}.) Say it once more.
\end{tcolorbox}
\section[\texorpdfstring{\ur{تپش} (tapish)}{تپش (tapish)}]{\ur{تپش} (tapish) — scorching heat}

\label{lesson:UR-C288-tapish}



\begin{quote}
Before the new one: say the Urdu for to glow red-hot, then the Urdu for brightness.
\end{quote}

\subsection*{You'll want to know: \ur{تپش}}
\textbf{\ur{تپش}} — \emph{tapish} — \textquotedblleft{}scorching heat\textquotedblright{}.

\begin{quote}
\ur{دوپہر} \ur{کی} \ur{تپش} \ur{سے} \ur{سب} \ur{تھک} \ur{گئے۔}
\end{quote}

\begin{quote}
\emph{dopahar kī tapish se sab thak gae.}
\end{quote}

\begin{quote}
Everyone was worn out by the scorching midday heat.
\end{quote}

Feminine.

\begin{cousinweb}
\textbf{\ur{تپش}} (\emph{tapish}) is usually traced to Persian \emph{tapish}, \textquotedblleft{}heat, a burning\textquotedblright{}. It is the heat you feel on the skin, of fire or of the sun.
\end{cousinweb}

\begin{grammarlens}[title={What you've built}]
One more word for telling what happened.
\end{grammarlens}

\subsection*{Your turn}
\begin{itemize}
\raggedright
  \item \emph{Say it:} \emph{tapish}
  \item \emph{Say it:} \emph{tapish}, once more
  \item \emph{Say it:} \emph{dopahar kī tapish se sab thak gae.} — Everyone was worn out by the scorching midday heat.
  \item \emph{Recall:} say the Urdu for to glow red-hot, then the Urdu for brightness, then say \emph{tapish} again
\end{itemize}

\begin{tcolorbox}[breakable,colback=teal!4,colframe=teal!35!black,title={Before you move on}]
What does \ur{تپش} mean? (\textquotedblleft{}Scorching heat\textquotedblright{}.) Say it once more.
\end{tcolorbox}
